\chapter{Tinjauan Pustaka dan Dasar Teori}

%% VERSI BAHASA INDONESIA. Konvensi peristilahan dijelaskan pada kepala chapter-1.tex:
%% prosa penghubung dalam bahasa Indonesia, istilah teknis dipertahankan dalam bahasa
%% Inggris (miring) kecuali padanannya benar-benar mapan.
%%
%% Setiap klaim mengenai sistem penulis lain telah diperiksa terhadap teks lengkap
%% sumbernya, dan penanda \fromabstract{} (dagger) sudah dilepas. Tiga klaim tidak lolos
%% pemeriksaan dan telah ditulis ulang: ke2025 (proses tanpa halusinasi hanya dicapai
%% GPT-4 atas panduan internasional, bukan kesepuluh model), degroot2024 (sampel "data EHR
%% nyata" hanya satu pasien), dan kraljevic2021 (8,8 miliar token adalah korpus pelatihan
%% self-supervised satu rumah sakit, bukan validasi tiga rumah sakit). Dua blok \crosscheck
%% (Subbab 2.1.2 dan 2.1.4) masih menandai temuan yang MEMBATASI klaim kebaruan pada Bab I
%% dan IV; zhou2022 dan ali2026 terbukti persis seperti yang dijelaskan, sehingga kedua bab
%% tersebut tetap perlu direvisi sebelum penyerahan.

\section{Tinjauan Pustaka}

Subbab ini menelaah penelitian terdahulu menurut tema dan menempatkan celah penelitian yang
dijawab skripsi ini.

\subsection{Pendukung keputusan berbasis panduan klinis}
Pendukung keputusan terkomputerisasi merupakan jalur yang sudah umum digunakan untuk mempersempit jarak antara panduan klinis dan praktiknya. Dampak dari adanya pendukung keputusan tersebut kini telah terukur dengan baik. Sebuah meta-analisis \cite{kwan2020} atas 122 uji klinis yang mencakup 1.203.053 pasien
menemukan bahwa pendukung keputusan meningkatkan proporsi pasien yang menerima perawatan
sesuai rekomendasi sebesar 5,8\% (IK~95\% 4,0---7,6), dengan heterogenitas yang tinggi ($I^2
= 76\%$). Perbaikan pada luaran klinisnya jauh lebih kecil, dengan median hanya 0,3\%. Nilai guna sistem semacam ini karenanya berdampak, tetapi masih sederhana dan terpusat pada
pengukuran proses alih-alih pengukuran luarannya. Pola yang sama juga muncul pada konteks
perioperatif. Di sana, penerapan panduan klinis membaik dan sebagian luaran keselamatan
turut membaik, tetapi mortalitas pascaoperasi tidak \cite{cai2024}.

Implementasi secara tradisional yang dominan adalah berupa \emph{computer-in\-ter\-pret\-able guideline} (CIG), yaitu
pengodean logika panduan ke dalam formalisme khusus---Arden Syntax, GLIF, PROforma, Asbru,
EON, dan kemudian GLARE serta SAGE---yang dieksekusi oleh sebuah mesin terhadap data pasien
\cite{declercq2004,declercq2008,oliveira2014}. Belum ada satu formalisme pun yang diberlakukan
sebagai standar \cite{declercq2008} dan Oliveira \textit{et al.}
mengidentifikasi interoperabilitas antara mesin eksekusi panduan dan model data klinis
sebagai hambatan yang terus berulang \cite{oliveira2014}---dan hambatan
itulah yang digarap skripsi ini, dengan mencocokkan konsep panduan terhadap data EHR berkode
melalui ruang konsep bersama.

Secara empiris, hambatan tersebut belum teratasi. Sebuah tinjauan sistematis atas pendukung
keputusan berbasis aturan menyaring 1.731 artikel dan hanya menemukan
15 artikel yang melaporkan evaluasi luaran klinisnya; di antaranya,
hanya dua yang terintegrasi dengan sistem EHR, enam menggunakan ontologi, dan
tiga menggunakan standar teknologi kesehatan yang diakui
\cite{papadopoulos2022}. Hambatan yang dilaporkan pada literatur ini konsisten:
ketertutupan aturan yang dikodekan, \emph{alert fatigue}, ketidakcocokan dengan alur kerja,
serta metodologi evaluasi yang lemah
\cite{papadopoulos2022,bosetti2025,silva2022}. Dari sini, terdapat dua implikasi bagi skripsi
ini. Pertama, integrasi dengan EHR masih jarang dilakukan, belum menjadi norma---inilah yang
menjadikan rancangan FHIR-native sebuah kontribusi, bukan sekadar detail implementasi. Kedua,
langkanya evaluasi yang ketat menjadikan rancangan keselarasan-terhadap-rujukan-independen
pada Bab~III sebagai kontribusi metodologis tersendiri.

\subsection{\emph{Knowledge graph} untuk pendukung keputusan klinis}
\label{sec:rw-kg}
\emph{Knowledge graph} memadukan pengetahuan klinis yang heterogen dan data pasien ke dalam struktur yang dapat dikueri. Literatur kesehatan mengenainya sudah cukup banyak sehingga
telah beberapa kali ditinjau \cite{abusalih2023,cui_h2023,budhdeo2023}. Sebuah
\emph{scoping review} atas 255 artikel melaporkan ukuran graf dengan median
46.983 simpul (IQR 6.415---460.948) dan 906.737
sisi (IQR 66.272---9.894.909) serta menyimpulkan bahwa
penerjemahan \emph{knowledge graph} ke dalam praktik klinis masih menjadi
tantangan \cite{budhdeo2023}. Graf yang dilaporkan pada Bab~IV berada di atas median
tersebut baik pada simpul maupun pada sisinya, yaitu pada bagian atas dari rentang interkuartil yang dilaporkan.

Sistem terdahulu terbagi menjadi dua kelompok menurut apa yang dihubungkan olehnya. Di satu sisi terdapat graf
\emph{pasien}: Xiao \textit{et al.} memaparkan repositori OMOP sebagai \emph{knowledge graph}
FHIR RDF dan memvalidasinya melalui kesesuaian kueri terhadap basis data sumber, tetapi
tidak mengodekan panduan klinis sehingga tidak menghasilkan rekomendasi
\cite{xiao2022}; Shang \textit{et al.} membangun graf berorientasi EHR untuk pertimbangan
kolaboratif antarinstitusi, diterapkan untuk identifikasi penyakit ginjal kronis yang
pernah terlewat dari penilaian klinisi \cite{shang2024}. Di sisi lain, terdapat graf \emph{pengetahuan}: Cui
\textit{et al.} membangun graf berbasis panduan klinis untuk obat kardiovaskular di Neo4j, tetapi
tidak melekatkan rekam medis pasien \cite{cui2025}. Alat bantu keputusan
berbasis FHIR yang lebih awal, seperti karya Semenov \textit{et al.}, memasangkan lapisan
interoperabilitas FHIR dengan basis pengetahuan berstruktur graf, tetapi
bekerja sebagai sistem aturan produksi berbasis panel laboratorium yang terbatas,
bukan berbasis pada rekam medis longitudinal \cite{semenov2018}. Seneviratne \textit{et al.} paling
dekat tendensinya dengan penelitian ini karena memodelkan provenans panduan klinis dan kohort
studi agar rekomendasi dapat membawa bukti pendukungnya \cite{seneviratne2023}.

Satu kelompok perkakas (\emph{tool}) yang sudah lazim juga perlu disebut di sini karena perkakas tersebut yang
mengakibatkan adanya keterbatasan dari apa yang dapat diklaim oleh skripsi ini. Pemetaan dari FHIR ke \emph{property graph}---
\emph{resource} menjadi simpul, rujukan menjadi sisi---sudah bersifat mekanis,
mapan, dan tersedia sebagai perkakas siap pakai: \emph{bundle viewer} berbasis peramban (\emph{web})
menampilkan sebuah \emph{bundle} unggahan sebagai graf simpul-sisi interaktif tanpa menuntut
satu pun keputusan pemodelan dari penggunanya \cite{tenasol_viewer}. Karena itu, transformasi FHIR ke \emph{property graph}
tidak diklaim sebagai kontribusi penelitian ini. Meski begitu, perkakas tersebut hanya bekerja terhadap satu \emph{bundle}
pada satu waktu dan tidak menyimpannya sehingga tidak ada pertanyaan lintas pasien yang dapat
diajukan. Perkakas hanya menampilkan setiap \emph{coding} secara \emph{inline} di dalam
\emph{resource} yang membawanya, alih-alih mengangkat \emph{coding} menjadi simpul bersama
sehingga tidak ada dua pasien yang tersambung melalui kode yang sama. Perkakas tersebut tidak
mencocokkan apa pun antara SNOMED~CT, RxNorm, dan LOINC. Perkakas juga tidak memuat
pengetahuan yang tidak berasal dari \emph{bundle} sehingga tidak ada lapisan panduan klinis maupun
literatur yang dapat dijadikan sebagai dasar penalaran. Oleh karena itu, sebuah \emph{viewer}---seperti perkakas tersebut---justru
dituntut \emph{faithful} terhadap masukannya sebab tujuannya adalah inspeksi. Sebaliknya,
graf yang dibangun pada Bab~III sengaja dibuat \emph{lossy}: graf tersebut membuang
\emph{resource} penagihan (\emph{billing}) dan provenans serta seluruh kode observasi LOINC
kecuali sembilan kode yang berada dalam \emph{allow-list}. Alasannya, graf itu dibangun
untuk menjawab satu pertanyaan klinis, bukan untuk merepresentasikan sebuah dokumen.

Ada satu sistem yang menggabungkan keduanya dan sistem itulah yang membuat klaim
kebaruan skripsi ini dinyatakan secara lebih sempit. Zhou \textit{et al.} membangun
\emph{knowledge graph} pengobatan hipertensi dari panduan hipertensi Tiongkok dan
mengevaluasinya terhadap 124 rekam medis yang dikumpulkan dari seorang dokter
spesialis di satu rumah sakit, dengan melaporkan \emph{recall} 91\%, hit@3
83\%, dan MRR 77\% pada tugas rekomendasi kelas obat tunggal \cite{zhou2022}. Dengan
demikian, lapisan panduan klinis, rekam medis pasien, dan evaluasi rekomendasi sudah berpadu dalam
satu sistem sehingga gabungan ketiganya \emph{tidak} diklaim sebagai kebaruan di sini. Penelitian ini berbeda dalam cakupan dan evaluasinya, bukan dalam arsitekturnya, dan jaraknya pun lebih sempit daripada yang dilaporkan sistem mereka. Graf mereka mencakup seluruh
kelas obat antihipertensi dan membawa \emph{compelling level} dari panduan klinis ke dalam
aturan penskorannya. Baik cakupan lintas kelas maupun kekuatan rekomendasi karenanya tidak
menjadi pembeda. Pembedanya terletak pada keluasan panduan klinis yang dikodekan: di sini
panduan itu mencakup antiplatelet, lipid, tekanan darah, penyekat beta,
renin-angiotensin-aldosteron, gaya hidup, serta kontraindikasi keamanan, alih-alih terapi
obat untuk satu kondisi saja. Nilai fisiologis
numerik pun turut masuk ke jalur penalaran mereka: atribut fisik merupakan satu dari empat
jenis entitas dalam graf mereka dan sebuah indikasi dapat dikualifikasi oleh ambang, seperti
pada contoh mereka sendiri berupa penyekat kanal kalsium dihidropiridin yang diindikasikan
untuk angina stabil dan hipertensi pada denyut jantung di bawah 75~bpm. Perbedaan keduanya terletak pada cara nilai tersebut digunakan---diketikkan dokter pada formulir entri pasien
untuk 124 rekam medis di sana, sedangkan di sini sebagai \emph{resource} FHIR yang membawa satuan
serta sistem kodenya sendiri, pada skala kohort dan tanpa penyalinan manual. Evaluasi mereka
melaporkan metrik pemeringkatan dari daftar kandidat yang sesuai, sedangkan keselarasan di sini diukur
terhadap standar rujukan yang independen dari sistem rekomendasi serta standar emas
berlabel manual. Terakhir, sistem mereka menyatukan dua sumber, sedangkan di sini lapisan
literatur disatukan sebagai sumber ketiga.

\subsection{Integrasi terminologi}
Pencocokan data pasien berkode terhadap konsep dari panduan klinis menuntut rekonsiliasi beberapa
kosakata---SNOMED~CT, RxNorm, dan LOINC---dan UMLS Metathesaurus merupakan cara baku
untuk menjembataninya \cite{umls}. Pemetaan ini tidak bebas dari kehilangan informasi: perbandingan
pemetaan ATC ke RxNorm menemukan 1.949 istilah tidak terpetakan pada UMLS dibanding
916 istilah pada OHDSI \emph{Standardized Vocabularies} meskipun UMLS mencapai
cakupan lebih tinggi (96,5\% dibanding 91\%) pada sampel data EHR milik satu pasien,
sebuah selisih yang bersumber dari dua kelas \emph{multi-ingredient} ATC \cite{degroot2024}.
Setiap rancangan jembatan konsep karenanya menanggung laju kegagalan pemetaan yang
seharusnya diukur, bukan diasumsikan.

Untuk menautkan teks bebas ke konsep, MedCAT menyediakan ekstraksi konsep UMLS dan
SNOMED~CT yang dilatih secara \emph{self-supervised} terhadap sekitar 8,8 miliar token rekam
medis satu rumah sakit dan tervalidasi di tiga rumah sakit di London, dengan keteralihan
yang dilaporkan sebesar $F_1 > 0{,}94$ antarlokasi \cite{kraljevic2021}. Penelitian ini
menggunakan scispaCy \cite{scispacy} atas dasar kesesuaian tugas, bukan kualitas. Keunggulan
MedCAT terletak pada \emph{disambiguation} berbasis konteks yang dipelajari secara
\emph{self-supervised} dari korpus narasi klinis. Kedua prasyarat itu tidak terpenuhi di
sini. Tugas penjembatanan pada Subbab~\ref{sec:kg-construction} bekerja pada kalimat
deskripsi dari kode tersebut---\texttt{Chronic ischemic heart disease (disorder)}---yang
tidak membawa kalimat pendukung lain untuk dibaca model semacam itu. Selain itu, ekspor
FHIR menyajikan data berkode alih-alih naratif sehingga tidak tersedia korpus untuk
melatihnya. Pustaka scispaCy juga memaparkan basis pengetahuannya sebagai indeks alias yang dapat
dikueri, bukan semata sebagai penganotasi atas teks. Sifat itulah yang memungkinkan
penelusuran alias secara persis serta penyaringan tipe semantik yang dipakai di sini.
Pustaka ini pun tidak menyimpan keadaan hasil pelatihan sehingga untai yang sama
menghasilkan konsep yang sama pada setiap proses. Skripsi ini bersandar pada sifat tersebut
ketika menyatakan bahwa penelusurannya deterministik. Untuk catatan klinis
berupa teks bebas, yang maknanya harus ditentukan dari konteks, MedCAT justru merupakan pilihan
yang lebih tepat.

\subsection{Model bahasa besar, \emph{retrieval-augmented generation}, dan \emph{faithfulness}}
Model bahasa besar menghasilkan teks klinis yang fasih, tetapi dapat berhalusinasi;
\emph{retrieval-augmented generation} (RAG) meredam ini dengan mensyaratkan pembangkitan hanya berdasar pada
bukti yang ditelusuri \cite{rag_lewis2020}. Sebuah tinjauan sistematis RAG di bidang kesehatan
menemukan beragam arsitektur yang digunakan, tetapi melaporkan ketiadaan
kerangka evaluasi yang terbakukan untuk aplikasi berbasis RAG dan bahwa sebagian besar studi
tidak membahas pertimbangan etis \cite{amugongo2025}. Ketiadaan itulah yang
menjadikan pemeriksaan \emph{faithfulness} secara mekanis pada Bab~III layak dilaporkan
sebagai kontribusi.

Sistem RAG khusus panduan klinis kini melaporkan hasil yang kuat. Sebuah sistem RAG atas panduan NICE
Britania Raya meningkatkan \emph{faithfulness}---yaitu apakah suatu jawaban
didukung oleh teks sumbernya---sebesar 64,7 poin persentase menjadi
99,5\% \cite{lewis2025}. Sebuah kerangka untuk interpretasi
panduan hepatologi meningkatkan akurasi dari 43\% menjadi 99\% melalui
pemformatan ulang panduan klinis dan rekayasa \emph{prompt} \cite{kresevic2024}. Selain itu,
evaluasi LLM-RAG pada sepuluh model dengan 3.234 respons melaporkan ketiadaan
halusinasi pada konfigurasi terbaiknya, yaitu GPT-4 pada panduan klinis internasional, yang
juga mengungguli jawaban yang dibuat manusia. Adapun laju halusinasi pada model lain
mencapai 48,6\% \cite{ke2025}.

Penelitian terkait juga memasangkan \emph{knowledge graph} dengan model bahasa untuk
kemampuan penjelasannya \cite{park2025} sehingga kombinasi graf dan LLM merupakan
area penelitian yang aktif, bukan area yang belum dijelajahi.

\subsection{Data klinis sintetis dan penelitian berbasis FHIR}
Synthea menyimulasikan rentang hidup pasien melalui model penyakit modular dan mengekspornya
sebagai \emph{bundle} HL7 FHIR sehingga memungkinkan penelitian informatika klinis yang
dapat direproduksi tanpa risiko privasi \cite{walonoski2017}. Synthea kemudian diperluas
dengan data pencitraan, genomik, dan catatan klinis yang ditautkan melalui FHIR
\cite{walonoski2022}. Validitasnya telah ditelaah secara langsung dengan
membandingkan ukuran mutu klinis yang dihitung pada kohort 1,2 juta pasien Massachusetts
terhadap angka yang dipublikasikan. Di samping itu, Chen \textit{et al.} menyimpulkan bahwa Synthea
andal dalam memodelkan demografi dan peluang layanan diberikan, tetapi
kemampuannya memodelkan luaran kesehatan yang heterogen setelah layanan
diberikan terbatas dan bahwa Synthea tidak memodelkan penyimpangan dalam
tata laksana \cite{chen2019}. Perlu dicatat, mereka melaporkan bahwa
tidak ada penduduk Synthea yang tekanan darahnya terkendali setelah diagnosis
hipertensi, sebanding dengan angka Massachusetts sebesar 74,5\% \cite{chen2019}.

\subsection{Pergeseran peristilahan dan kerangka dari waktu ke waktu}
Penelusuran literatur ini membutuhkan dua kosakata. Pertama, karya yang terbit sebelum kurang lebih 2015
menggambarkan arsitektur yang sama dengan istilah \emph{Semantic Web}, yaitu ontologi OWL,
\emph{description-logic reasoner}, tripel RDF, dan \emph{proof engine}. Contohnya adalah
pendukung keputusan praoperatif berbasis ontologi milik Bouamrane \textit{et al.}
\cite{bouamrane2011} dan kerangka eksekusi panduan klinis milik Hussain
\textit{et al.} \cite{hussain2007}. Kerangka yang disebut terakhir bahkan telah memadukan
ontologi domain, ontologi panduan klinis, dan ontologi pasien serta membangkitkan
\emph{justification tree} untuk memaparkan jejak inferensi. Kedua, karya pasca-2020 menggambarkan
persoalan yang sama dengan kosakata \emph{property graph}---\emph{knowledge graph}, Neo4j,
Cypher, GraphRAG. Arsitektur yang diusulkan skripsi ini karenanya merupakan wujud kekinian dari
rancangan yang telah lama mapan dan apabila penelusuran dibatasi hanya pada istilah ``\emph{knowledge
graph}'' akan berakibat terlewatnya satu dasawarsa preseden yang relevan secara langsung.

\subsection{Ringkasan dan celah penelitian}
Penelitian terdahulu sebagian besar menggarap tema-tema ini secara terpisah. Sistem berbasis
aturan dan CIG mengodekan logika panduan klinis, tetapi jarang terintegrasi dengan EHR berstandar
\cite{papadopoulos2022}. Graf klinis sisi pasien pernah dibangun secara ketat, tetapi
tidak membawa lapisan panduan klinis \cite{xiao2022}. Graf sisi panduan klinis tidak membawa rekam medis
pasien \cite{cui2025}. Satu sistem memang membawa keduanya, yaitu graf obat berbasis panduan klinis
yang dievaluasi terhadap rekam medis nyata \cite{zhou2022}, tetapi terbatas pada terapi obat
satu kondisi dan dievaluasi dengan metrik pemeringkatan alih-alih terhadap standar rujukan.
Model bahasa yang dilandaskan pada graf terbukti menurunkan halusinasi
\cite{ali2026}, tetapi dievaluasi pada akurasi penjawaban pertanyaan---enam puluh
pertanyaan dalam satu area penyakit, dinilai terhadap jawaban rujukan yang ditarik dari
literatur---bukan pada keselarasan terhadap panduan klinis terhadap suatu kohort pasien. Kontribusi skripsi ini karenanya lebih sempit dan lebih tepat daripada sekadar ``menyatukan
data pasien dengan panduan klinis''. Pertama, skripsi ini membangun graf FHIR-native yang di
dalamnya \emph{observasi terstruktur} tersedia untuk menggerbangi rekomendasi. Kedua, graf
itu dievaluasi berdasarkan keselarasan panduan klinis terhadap standar rujukan yang
independen dari sistem rekomendasi. Ketiga, graf itu disertai \emph{grounded briefing} yang
\emph{faithfulness}-nya pada tingkat sitasi dapat diperiksa secara mekanis karena
penelusurannya deterministik dan membawa provenans.

\section{Dasar Teori}

\subsection{Penyakit jantung koroner kronis dan panduan AHA/ACC 2023}
Penyakit jantung koroner kronis (CCD) merupakan fase stabil dari penyakit arteri koroner yang
memerlukan pencegahan sekunder jangka panjang.

\subsubsection*{Catatan peristilahan: SIHD dan CCD}
Kedua istilah mengenai jantung koroner kronis yang digunakan dalam skripsi ini berkaitan, tetapi tidak dapat saling menggantikan dan
pembedaannya penting karena data EHR dan panduan klinis berasal dari sisi yang berbeda. Hingga 2023,
kondisi ini dibahas dengan sebutan \emph{stable ischemic heart disease} (SIHD), yaitu pasien yang
penyakit koronernya telah terdiagnosis namun tidak sedang berada dalam fase akut. Panduan AHA/ACC
2023 \cite{virani2023} menggantikan kerangka tersebut dengan \emph{chronic coronary disease},
sebuah entitas yang sengaja dibuat lebih luas: selain angina stabil, panduan itu secara
eksplisit mencakup pasien yang telah lebih dari satu tahun melewati sindrom koroner akut,
pasien pascarevaskularisasi, pasien dengan kardiomiopati iskemik, dan pasien yang penyakit
koronernya terdokumentasi tanpa gejala saat ini. SIHD karenanya paling tepat dibaca sebagai
\emph{subhimpunan} dari CCD---setiap pasien SIHD termasuk CCD, tetapi tidak setiap pasien CCD dapat
disebut SIHD.

Ketidaksetangkupan ini tampak dalam penelitian ini. Korpus sintetis dibangkitkan dengan modul
Synthea bernama \texttt{stable\_ischemic\_heart\_disease}, yang kosakatanya mendahului
penggantian nama pada 2023, sedangkan panduan klinis yang dikodekan adalah panduan CCD. Modul
tersebut karenanya memasok pasien dari kategori lama yang lebih sempit, sementara
rekomendasinya ditulis untuk kategori yang lebih luas. Akibatnya, terdapat dua hal. Keanggotaan
kohort tidak bergantung pada modul pembangkit mana yang digunakan, melainkan
oleh diagnosis berkode yang benar-benar ada dalam rekam medis (Subbab~\ref{sec:gate})
sehingga pasien dari proses pembangkitan populasi umum yang mengalami diagnosis penentu tetap
disertakan atas dasar yang sama. Selain itu, karena CCD merupakan kategori yang lebih luas,
rekomendasi yang ditujukan pada bagian CCD di luar gambaran SIHD klasik---misalnya yang
dikaitkan dengan sindrom koroner akut terkini---lebih jarang terpakai oleh korpus ini
dibandingkan pada populasi CCD nyata. Panduan AHA/ACC CCD 2023 \cite{virani2023}
menata rekomendasinya menurut \emph{Class of Recommendation} (COR) dan
\emph{Level of Evidence}. COR terbagi menjadi Kelas~1 (kuat), Kelas~2a (sedang),
Kelas~2b (lemah), serta Kelas~3 yang terbagi lagi menjadi 3-No-Benefit dan 3-Harm, yakni
kontraindikasi. Dalam penelitian ini,
COR menentukan semantik sisi rekomendasi: rekomendasi Kelas~3 menjadi kontraindikasi,
sedangkan rekomendasi Kelas~1/2a/2b menjadi rekomendasi positif.

\subsection{HL7 FHIR}
\emph{Fast Healthcare Interoperability Resources} (FHIR) \cite{fhir_r4} adalah standar HL7
untuk pertukaran data kesehatan dalam bentuk \emph{resource} modular (misalnya
\texttt{Patient}, \texttt{Condition}, \texttt{Observation}, \texttt{MedicationRequest}, dan
\texttt{Procedure}). Antar-\emph{resource} saling merujuk dan makna klinis dibawa oleh elemen
\texttt{CodeableConcept}/\texttt{Coding} yang masing-masing memasangkan URI \texttt{system}
kode dengan sebuah \texttt{code}. Yang menentukan bagi penelitian ini, hasil numerik merupakan
\emph{resource} \texttt{Observation} yang memiliki \texttt{valueQuantity} sehingga
nilai laboratorium dan tanda vital tersedia sebagai data terstruktur. Rilis HL7 FHIR yang digunakan di
sepanjang skripsi ini adalah FHIR~R4; Subbab~\ref{sec:data-source} mencatat bagaimana hal itu
dipastikan dari instans yang diekspor.

Karena bagian selanjutnya dari skripsi ini sangat bergantung pada struktur tersebut, sebuah
contoh ringkas diberikan pada Listing~\ref{lst:fhir}. Contoh itu diambil apa adanya dari
sebuah \emph{bundle} hasil pembangkitan, dengan elemen yang tidak terpakai dihilangkan.

\begin{lstlisting}[caption={Cuplikan ringkas \emph{resource} FHIR R4 dari \emph{bundle}
pasien hasil pembangkitan: sebuah diagnosis berkode dan sebuah hasil laboratorium numerik.},
label={lst:fhir},
basicstyle=\ttfamily\scriptsize,breaklines=true,frame=single,captionpos=b]
{ "resourceType": "Condition",
  "id": "22f599b6-...-8e66-1f64bdc5c863",
  "code": {                              // APA temuannya
    "coding": [ { "system":  "http://snomed.info/sct",
                  "code":    "414545008",
                  "display": "Ischemic heart disease (disorder)" } ],
    "text": "Ischemic heart disease (disorder)" },
  "subject": { "reference": "urn:uuid:22f599b6-...-e586-..." },   // MILIK SIAPA
  "onsetDateTime": "2013-05-24T09:14:31+07:00" }

{ "resourceType": "Observation",
  "status": "final",
  "code": {                              // PENGUKURAN yang mana
    "coding": [ { "system":  "http://loinc.org",
                  "code":    "18262-6",
                  "display": "Cholesterol in LDL [Mass/volume] ..." } ] },
  "subject": { "reference": "urn:uuid:22f599b6-...-e586-..." },
  "valueQuantity": {                     // ANGKANYA -- dapat dibandingkan mesin
    "value": 91.35, "unit": "mg/dL",
    "system": "http://unitsofmeasure.org", "code": "mg/dL" } }
\end{lstlisting}

Tiga sifat dari struktur ini penting bagi keseluruhan sistem. Pertama,
antar-\emph{resource} \emph{ditautkan melalui rujukan}: sebuah \texttt{Condition} tidak
memuat pasiennya, melainkan menunjuk kepadanya---inilah sebabnya model graf pada Bab~III
mengubah rujukan menjadi sisi. Kedua, makna klinis dibawa \emph{bukan oleh nama elemen, melainkan oleh kodenya}. Elemen
\texttt{code} yang sama memuat konsep SNOMED~CT di sini dan akan memuat konsep RxNorm atau
LOINC di tempat lain. Itulah sebabnya setiap \emph{coding} diangkat ke satu simpul bersama
dan dijembatani ke ruang konsep yang sama. Ketiga, dan
menentukan bagi skripsi ini, sebuah hasil laboratorium bukanlah kalimat, melainkan
\texttt{valueQuantity} yang memuat angka, satuan, dan sistem satuan. Ekspor tabular
terdatarkan tetap membawa angka-angka tersebut, tetapi mengubah struktur yang menyertainya:
satu pembacaan berkomponen jamak terpecah ke beberapa baris yang harus dirakit ulang sebelum
sebuah ambang batas dapat dievaluasi. \emph{Resource} FHIR tetap mempertahankan nilai, satuan, dan keanggotaan panelnya. Sifat itulah yang memungkinkan sebuah rekomendasi yang bergantung pada ambang
batas diputuskan, bukan sekadar rekomendasi yang bergantung pada ada tidaknya suatu diagnosis
(Subbab~\ref{sec:flattening}).

\subsection{Terminologi klinis dan UMLS}
Tipe \emph{resource} yang berbeda dikodekan dalam kosakata yang berbeda: SNOMED~CT untuk
kondisi temuan dan prosedur klinis \cite{snomedct}, RxNorm untuk obat \cite{rxnorm}, dan LOINC untuk
observasi laboratorium dan tanda vital \cite{loinc}. \emph{Unified Medical Language System}
(UMLS) Metathesaurus \cite{umls} memadukan kosakata-kosakata tersebut beserta banyak lainnya
dan memberikan setiap konsep sebuah \emph{Concept Unique Identifier} (CUI), yang menyediakan
ruang konsep bersama tempat setiap kode EHR dan setiap target dari panduan klinis dipetakan dalam skripsi
ini.

\subsection{\emph{Knowledge graph} dan basis data \emph{property graph}}
\emph{Knowledge graph} merepresentasikan entitas sebagai simpul dan hubungan antarentitas
sebagai sisi berlabel. Pada basis data \emph{property graph} seperti Neo4j \cite{neo4j},
simpul dan sisi dapat membawa properti pasangan kunci-nilai (\emph{key-value}); data dikueri dengan bahasa
deklaratif Cypher \cite{cypher} melalui pencocokan pola. \emph{Property graph} sangat sesuai
untuk menyatukan sumber yang heterogen karena tipe simpul dan sisi baru dapat ditambahkan
secara bertahap dan hubungan lintas beberapa lompatan (\emph{multi-hop}) dapat ditelusuri secara langsung.

\subsection{Pemrosesan bahasa alami klinis: NER dan \emph{entity linking}}
Ekstraksi informasi terstruktur dari teks biomedis tak terstruktur memerlukan
\emph{named-entity recognition} (NER), yang mengenali segmen teks bermakna klinis, dan
\emph{entity linking}, yang memetakan segmen tersebut ke suatu basis pengetahuan seperti
UMLS. Pustaka scispaCy \cite{scispacy} menyediakan model NLP biomedis serta \emph{entity
linker} UMLS/MeSH yang digunakan dalam penelitian ini untuk memproses korpus literatur dan
membantu penjembatanan terminologi; model yang dipakai di sini adalah
\texttt{en\_core\_sci\_md}, yang vektornya dilatih atas OntoNotes~5, Common Crawl, dan
GENIA~1.0.

\subsection{Kekuatan asosiasi: \emph{pointwise mutual information}}
\emph{Pointwise mutual information} (PMI) mengukur seberapa jauh dua peristiwa---di sini
sebuah obat dan sebuah kondisi---muncul bersama lebih sering daripada yang diperkirakan jika
keduanya saling bebas \cite{pmi}. Untuk dua konsep $x$ dan $y$, PMI didefinisikan sebagai

\begin{equation}
\mathit{PMI}(x, y) \;=\; \log \frac{P(x, y)}{P(x)\,P(y)} .
\label{eq:pmi}
\end{equation}

Peluangnya dihasilkan dari frekuensi dokumen pada korpus abstrak. Dengan $N$ menyatakan
jumlah abstrak, $\mathit{df}(x)$ jumlah abstrak yang menyebut konsep $x$, dan $c(x,y)$ jumlah
abstrak yang menyebut keduanya, estimator yang disubstitusikan ke Persamaan~\ref{eq:pmi} adalah

\begin{equation}
\mathit{PMI}(x, y)
\;=\; \log_{10} \frac{c(x,y)/N}{\bigl(\mathit{df}(x)/N\bigr)\bigl(\mathit{df}(y)/N\bigr)}
\;=\; \log_{10} \frac{c(x,y)\;N}{\mathit{df}(x)\,\mathit{df}(y)} .
\label{eq:pmi-df}
\end{equation}

Basis~10 digunakan sehingga nilai $1$ menandakan pasangan yang muncul bersama sepuluh kali
lebih sering daripada perkiraan kebebasan dan nilai $0$ menandakan kebebasan secara penuh;
nilai negatif menandakan kemunculan bersama yang lebih jarang daripada peluang acak. Bentuk
ruas kanan Persamaan~\ref{eq:pmi-df} adalah bentuk yang dihitung.

PMI digunakan untuk memeringkat asosiasi antara obat dan kondisi yang ditambang dari literatur. PMI
diterapkan dengan ambang minimum jumlah artikel: karena $c(x,y)$ berada pada pembilang, sebuah
pasangan yang hanya muncul pada satu abstrak dapat memperoleh skor tinggi semata-mata karena
kedua konsepnya langka sehingga pasangan di bawah ambang tersebut dibuang alih-alih
masuk dalam proses pemeringkatan. Ambangnya adalah tiga artikel: sebuah pasangan obat--kondisi harus disebut
bersama dalam sekurang-kurangnya tiga abstrak berbeda sebelum diperingkatkan. Ambang
tersebut diterapkan saat \emph{briefing} dirakit, bukan saat sisi \emph{co-occurence} dibangun,
sehingga lapisan literatur tetap menyimpan setiap pasangan yang teramati dan ambangnya dapat
diubah tanpa mewujudkan ulang graf.

\subsection{\emph{Retrieval-augmented generation} dan peringkasan membumi}
\emph{Retrieval-augmented generation} (RAG) \cite{rag_lewis2020} mengondisikan keluaran model
bahasa terhadap konteks yang ditelusuri; RAG klasik menelusuri potongan teks tak terstruktur
melalui kemiripan vektor (\emph{dense}) sehingga menghasilkan konteks yang
\emph{aproksimatif}. Skripsi ini sebaliknya menggunakan penelusur \emph{knowledge graph} yang
deterministik menggunakan Cypher sehingga fakta yang ditelusuri bersifat persis dan membawa provenans. Model bahasa hanya melakukan \emph{grounded summarisation}---merumuskan ulang fakta tanpa
menambahkan informasi. \emph{Faithfulness} menyatakan sifat bahwa setiap pernyataan dalam
keluaran didukung oleh fakta yang sudah ditelusuri; pelanggarannya disebut halusinasi.

\subsection{Pengukuran kesepakatan: Cohen's kappa}
Cohen's $\kappa$ \cite{cohen1960} mengukur kesepakatan antara dua penilai---atau antara sebuah
metode otomatis dan standar emas---dengan koreksi terhadap kesepakatan yang diharapkan terjadi
secara kebetulan. Nilai di atas~0,8 secara konvensional ditafsirkan sebagai kesepakatan yang
hampir sempurna, sedangkan nilai 0,21---0,40 ditafsirkan sebagai kesepakatan yang cukup
(\emph{fair}) \cite{landis1977}. Nilai $\kappa$ juga dipengaruhi prevalensi kategori dan bias
antarpenilai sehingga indikator keduanya sebaiknya dilaporkan bersama $\kappa$
\cite{byrt1993}. Statistik ini digunakan dalam penelitian ini untuk memvalidasi standar
rujukan otomatis terhadap label manual dan untuk mengukur kesepakatan antarklinisi pada
telaah ketepatan rekomendasi.

\subsection{Data pasien sintetis: Synthea}
Synthea \cite{synthea} adalah pembangkit sumber terbuka untuk rekam medis pasien yang sintetis
namun realistis secara statistik, dihasilkan dengan menyimulasikan perjalanan penyakit melalui
model modular dan dapat diekspor sebagai \emph{bundle} HL7 FHIR maupun berkas CSV. Karena
datanya sintetis, data tersebut tidak memuat informasi kesehatan yang perlu dilindungi dan mendukung
penelitian yang sepenuhnya dapat direproduksi.

\section{Analisis Perbandingan Metode}

Pilihan rancangan dalam skripsi ini dijustifikasi melalui perbandingan dengan alternatifnya.
CDSS yang murni \emph{berbasis aturan} dapat mengodekan logika panduan klinis, tetapi buruk dalam
berintegrasi dengan data EHR berkode yang heterogen dan rapuh untuk dipelihara; pendekatan
\emph{knowledge graph} dengan penjembatanan UMLS sebaliknya mencocokkan data pasien berkode
terhadap konsep panduan klinis melalui ruang konsep bersama dan menelusuri hubungan lintas sumber
secara langsung. Untuk keluaran bahasa alami, peringkas \emph{murni-LLM} berisiko
berhalusinasi dan peringkas \emph{RAG-vektor} hanya berlandas pada potongan teks yang
ditelusuri secara aproksimatif; \emph{penelusuran knowledge graph yang deterministik} diikuti
perumusan kalimat oleh LLM menghasilkan fakta yang persis dan membawa provenans. Karena itu,
keluarannya dapat diperiksa terhadap himpunan fakta yang tertutup, hal yang esensial pada
konteks klinis. Pemeriksaan semacam itu belum mencakup kesetiaan makna setiap kalimat.
Terakhir, evaluasi pada observasi terstruktur yang \emph{FHIR-native}
mempertahankan setiap nilai tetap terikat pada satuan, sistem satuan, dan panelnya sehingga
rekomendasi bergantung-observasi dapat digerbangi menurut nilai terukur tanpa perlu merakit
ulangnya dari ekspor tabular terdatarkan---mekanisme yang dapat meningkatkan presisi.
Tabel~\ref{tab:method-comparison} merangkum perbandingan tersebut.

\begin{table}[h]
\centering
\caption{Alternatif rancangan sistem pendukung keputusan klinis untuk CCD, dibandingkan
terhadap kriteria yang melatarbelakangi skripsi ini.}
\label{tab:method-comparison}
\footnotesize
\setlength{\tabcolsep}{3pt}
\begin{tabular}{p{2.0cm}p{2.28cm}p{2.28cm}p{2.28cm}p{2.28cm}p{2.28cm}}
\hline
\textbf{Pendekatan} & \textbf{Integrasi EHR berkode} & \textbf{Penjembatanan terminologi} &
\textbf{Penggerbangan observasi numerik} & \textbf{Penyatuan lintas sumber} &
\textbf{\emph{Faithfulness} keluaran} \\
\hline
CDSS berbasis aturan/CIG &
Jarang dalam praktik: 2 dari 15 sistem terevaluasi \cite{papadopoulos2022} &
Dikodekan manual per aturan &
Dapat dinyatakan, tetapi manual per aturan &
Hanya panduan &
Tidak berlaku (tanpa keluaran bahasa alami) \\
Basis data relasional dengan SQL &
Didukung langsung &
Perlu tabel pemetaan eksplisit &
Didukung langsung &
Perlu skema dan migrasi baru per sumber &
Tidak berlaku (tanpa keluaran bahasa alami) \\
\emph{RAG-vektor} atas teks panduan &
Hanya teks; kode tidak digunakan &
Kemiripan \emph{embedding} &
Tidak terepresentasi &
Tingkat korpus, bukan entitas &
Aproksimatif; berlandas potongan teks tertelusur \\
\emph{Prompting} murni-LLM &
Hanya teks \emph{prompt} &
Implisit dan tak terverifikasi &
Tidak terepresentasi &
Tidak ada &
Tanpa kendala; berisiko halusinasi \\
\textbf{Skripsi ini} &
\emph{FHIR-native}, terstruktur &
Ruang konsep bersama UMLS &
Digerbangi nilai terukur &
EHR, panduan, dan literatur dalam satu graf &
Penelusuran deterministik; terverifikasi otomatis \\
\hline
\end{tabular}
\end{table}

Salah satu baris pada Tabel~\ref{tab:method-comparison} merupakan
alternatif terdekat dengan rancangan yang dipakai di sini. Basis data relasional atas data FHIR yang sama juga akan menyimpan observasi numerik secara
langsung. Lintasan inferensi kanonis yang sama pun dapat dinyatakan di sana dalam bentuk
\emph{join} berganda yang setara. Karena itu, pemilihan graf properti merupakan soal
\emph{kesesuaian representasi}, bukan soal kemampuan secara mentah. Skripsi ini tidak
mengklaim bahwa alternatif relasional tidak mampu melakukannya. Meski
demikian, tiga pertimbangan mendukung penggunaan graf. Pertama, FHIR mendefinisikan lebih dari seratus jenis \emph{resource} \cite{fhir_r4}. Skema
relasional yang \emph{faithful} memerlukan satu tabel per jenis serta satu migrasi untuk
setiap jenis baru yang dimuat. Sebaliknya, rancangan label ganda
\texttt{:Resource:}\textit{Type} pada Bab~III menerima jenis \emph{resource} apa pun melalui
satu jalur pemuatan tunggal. Kedua, ketiga sumber yang disatukan, yaitu rekam medis pasien,
sekumpulan rekomendasi bersyarat, dan jaringan kemunculan bersama literatur, memiliki
sedikit kesamaan struktur. Lapisan tambahan pun dapat dipasang pada graf tanpa mengganggu
lapisan yang telah ada. Ketiga, dan paling relevan bagi RQ3, penelusuran yang ditulis sebagai rangkaian
hubungan klinis dapat diaudit sebagai rangkaian penalaran klinis---hal yang penting bagi sistem
yang keluarannya harus membawa provenans.

\section{Pertanyaan Tugas Akhir}

Setiap masalah yang dirumuskan pada Bab~I diubah menjadi pertanyaan yang dapat dijawab
secara empiris. Pertanyaan penelitian berikut karenanya berpadanan satu lawan satu dengan
masalah tersebut dan dipetakan ke tujuan penelitian:

\begin{enumerate}
    \item[RQ1.] Bagaimana \emph{knowledge graph} FHIR-native terpadu yang menyatukan EHR,
    panduan klinis, dan literatur dapat dibangun untuk CCD, dengan terminologi klinis dijembatani ke
    dalam satu ruang konsep bersama? (memetakan ke RO1)
    \item[RQ2.] Apakah pemuatan observasi terstruktur meningkatkan keselarasan panduan klinis pada
    sistem rekomendasi dibandingkan \emph{baseline} yang hanya memanfaatkan
    kondisi/teks saja? (memetakan ke RO2)
    \item[RQ3.] Seberapa \emph{faithful} \emph{grounded briefing} yang dihasilkan LLM dan
    sejauh mana penelusuran \emph{knowledge graph} yang deterministik dapat menjamin
    \emph{faithfulness} tersebut? (memetakan ke RO3)
\end{enumerate}
